% from ega0-20-fr.tex
\oldpage[0\textsubscript{IV-1}]{127}
\begin{corollary}[20.4.10]
\label{0.20.4.10-fr}
Supposons en outre que dans $B$ le carré de tout idéal ouvert soit ouvert ;
alors, pour tout idéal ouvert $\mathfrak{R}$ de $B$,
$\Omega^1_{B/A}\otimes_B(B/\mathfrak{R})
=\Omega^1_{B/A}/\mathfrak{R}.\Omega^1_{B/A}$ est un
$(B/\mathfrak{R})$-module projectif.
\end{corollary}

\begin{proof}
En effet, la topologie de $\Omega^1_{B/A}$ est alors déduite de celle de
$B$ \sref{0.20.4.5-fr}, et il suffit d'appliquer \sref{0.19.2.4-fr}.
\end{proof}

\begin{corollary}[20.4.11]
\label{0.20.4.11-fr}
Soient $A$, $B$ deux anneaux locaux noethériens, $\rho:A\to B$ un
homomorphisme local, faisant de $B$ une $A$-algèbre topologique formellement
lisse (pour les topologies préadiques). Alors, pour tout idéal de définition
$\mathfrak{b}$ de $B$,
$\Omega^1_{B/A}\otimes_B(B/\mathfrak{b})$ est un
$(B/\mathfrak{b})$-module libre.
\end{corollary}

\begin{proof}
En effet, il résulte de \sref{0.20.4.10-fr} que ce module est projectif, et
comme $B/\mathfrak{b}$ est un anneau artinien, tout
$(B/\mathfrak{b})$-module projectif est libre
($\mathbf{0_{III}}$, 10.1.3).
\end{proof}

\begin{proposition}[20.4.12]
\label{0.20.4.12-fr}
Soient $A$ un anneau topologique, $B$ une $A$-algèbre topologique. Si
l'homomorphisme structural $A\to B$ est surjectif, on a
$\Omega^1_{B/A}=0$.
\end{proposition}

\begin{proof}
En effet, on a $\operatorname{D\acute{e}r}_A(B,L)=0$ pour tout $B$-module
$L$ en vertu de \sref{0.20.1.1-fr}, et la proposition découle donc aussitôt
de \sref{0.20.4.8-fr}.
\end{proof}

\begin{examples}[20.4.13]
\label{0.20.4.13-fr}
\begin{enumerate}
  \item[\rm(i)] Soient $A$ un anneau,
  $B=A[X_\alpha]_{\alpha\in I}$ une algèbre de polynômes sur $A$. Alors
  $\Omega^1_{B/A}$ est un $B$-module libre, dont les $dX_\alpha$ forment une
  base. En effet, les $dX_\alpha$ engendrent ce $B$-module
  \sref{0.20.4.7-fr}. D'autre part, si $L$ est un $B$-module libre, ayant
  une base $(z_\alpha)_{\alpha\in I}$ dont $I$ est l'ensemble d'indices, il
  existe un $A$-homomorphisme $u$ de $B$ dans $D_B(L)$ tel que
  $u(X_\alpha)=(X_\alpha,z_\alpha)$, pour tout $\alpha\in I$, donc
  \sref{0.20.1.5-fr} une $A$-dérivation $D$ de $B$ dans $L$ telle que
  $D(X_\alpha)=z_\alpha$ pour tout $\alpha$ ; en vertu de
  \sref{0.20.4.8.1-fr}, cela prouve que les $dX_\alpha$ sont linéairement
  indépendants.
  \item[\rm(ii)] Soient $A$ un anneau, $L$ un $A$-module, $B$ la
  $A$-algèbre $D_A(L)$ ; alors l'homomorphisme canonique
  $x\mapsto d_{B/A}x$ de $L$ dans $\Omega^1_{B/A}$ est \emph{bijectif}, car
  il est immédiat que les $B$-dérivations de $B=D_A(L)$ dans un $B$-module
  $M$ sont les applications de la forme $(b,x)\mapsto u(x)$, où
  $u\in\operatorname{Hom}_A(L,M)$ ; on conclut donc par
  \sref{0.20.4.8-fr}, (ii).
  \item[\rm(iii)] Supposons que $B$ soit le produit de deux $A$-algèbres
  topologiques $B_1$, $B_2$ (identifiées à des idéaux de $B$). Alors les
  images de $B_1\otimes_A B_2$ et de $B_2\otimes_A B_1$ par l'homomorphisme
  $p$ \sref{0.20.4.1.1-fr} sont nulles, d'où il résulte aussitôt que
  $\mathrm{P}^1_{B/A}$ s'identifie au \emph{produit}
  $\mathrm{P}^1_{B_1/A}\times\mathrm{P}^1_{B_2/A}$, et que le $B$-module
  $\Omega^1_{B/A}$ est \emph{somme directe} (topologique) des $B$-modules
  $\Omega^1_{B_1/A}$ et $\Omega^1_{B_2/A}$ (annulés respectivement par
  $B_2$ et $B_1$).
  \item[\rm(iv)] Soient $A$, $B$ deux anneaux intègres tels que $A\subset
  B$, que $A$ soit intégralement clos, $B$ entier sur $A$, et que le corps
  des fractions de $B$ soit une extension séparable de celui de $A$. Alors
  le $B$-module $\Omega^1_{B/A}$ est un \emph{module de torsion}. En effet,
  pour tout $x\in B$, le polynôme minimal de $x$ par rapport au corps des
  fractions $K$ de $A$ est un polynôme $f(T)$ appartenant à $A[T]$ ; comme
  $x$ est séparable sur $K$, on a $f'(x)\neq0$ et d'autre part on déduit de
  la relation $f(x)=0$ que $f'(x)d_{B/A}x=0$, d'où notre assertion en vertu
  de \sref{0.20.4.7-fr}.
\end{enumerate}
\end{examples}

\begin{remarks}[20.4.14]
\label{0.20.4.14-fr}
\begin{enumerate}
  \item[\rm(i)] On notera que le $B$-module $\Omega^1_{B/A}$, privé de sa
  topologie, est indépendant des topologies de $A$ et de $B$.
  \item[\rm(ii)] Nous introduirons plus tard l'« algèbre des parties
  principales d'ordre $n$ » de $B$ par rapport à $A$,
  $\mathrm{P}^n_{B/A}=(B\otimes_A B)/(\mathfrak{J}_{B/A})^{n+1}$, qui est à
  la base du « calcul différentiel d'ordre $n$ ».
\end{enumerate}
\end{remarks}

\subsection{Propriétés fonctorielles fondamentales de
\texorpdfstring{$\Omega^1_{B/A}$}{Omega1(B/A)}}
\label{subsection:0.20.5-fr}

